%============================================================
% Chapter 30: Scanline over the Right Endpoint
% STATUS: PLANNED  (generated by tools/outline.py; replace this file when the chapter is written)
%============================================================
\chapter{Scanline over the Right Endpoint}
\label{ch:scanline}

\begin{tcolorbox}[enhanced, colback=white, colframe=secondaryblue, boxrule=1pt, arc=2mm]
\textbf{\large\textcolor{primaryblue}{For every right endpoint, keep the answer for all left endpoints at once --- the standard way to count or sum over all subarrays with a property.}}
\end{tcolorbox}

\chapterfacts{1900}{2600}{5}{10}{30}

\begin{plannedbox}[Planned Contents --- this chapter is not written yet]
\textbf{Part A --- Lessons}
\begin{enumerate}[leftmargin=*, label=\textbf{A\arabic*.}, itemsep=1pt]
    \item The scanline invariant: an array indexed by the left endpoint
    \item Range updates when the right endpoint moves (data structures are the tree book)
    \item Last occurrences and distinct-value scanlines
    \item Offline queries $[l, r]$ answered when the sweep reaches $r$
    \item Summing answers over all subarrays with historic sums
\end{enumerate}
\textbf{Part B --- Patterns} (with candidate problems, checked when the chapter is written)
\begin{enumerate}[leftmargin=*, label=\textbf{P\arabic*.}, itemsep=1pt]
    \item \textbf{Distinct-value scanline}
    \item \textbf{Range updates per step}
    \item \textbf{Queries at their right end}
    \item \textbf{Historic sums}
    \item \textbf{Two-dimensional scanlines}
\end{enumerate}
\end{plannedbox}
